\begin{lemma}\label{lem:trace}
 Let $k\geq 1$. We have: 
 $$
\dim \Hom_G(\rho^{\otimes k}, \C[G]^{adj}) = \sum_{g \in G//G} \chi_\rho(g)^k = tr(A^k_{\Gamma}) = \abs{\{\text{circuits of length } k \text{ in } \Gamma \}}.  $$

\end{lemma}
\begin{proof}
 The eigenvalues of the matrix $A_{\Gamma}$ are precisely $\chi_\rho(g)$ for $g \in G//G$ (see Proposition \ref{prop:spectra_adj_matrix}), so the sum of eigenvalues of $A^k_{\Gamma}$ is $\sum_{g \in G//G} \chi_\rho(g)^k $, which proves the middle equality. The right equality holds for any graph $\Gamma$ and is obvious.
 
 It remains to prove that $\dim \Hom_G(\rho^{\otimes k}, \C[G]^{adj}) = tr(A^k_{\Gamma})$. Indeed, we have:
 \begin{align*}
 \dim \Hom_G(\rho^{\otimes k}, \C[G]^{adj}) &= \dim \Hom_G(\rho^{\otimes k}, \bigoplus_{\mu \in \mathrm{Irr}(G)} \mu^* \otimes \mu) \\
 &= \sum_{\mu \in \mathrm{Irr}(G)} \dim \Hom_G(\rho^{\otimes k}\otimes \mu, \mu)
 \end{align*}
 In the right hand side, each summand $\dim \Hom_G(\rho^{\otimes k}\otimes \mu, \mu)$ is precisely the diagonal entry in $A^k_{\Gamma}$ corresponding to $\mu \in \mathrm{Irr}(G)$, so $$\sum_{\mu \in \mathrm{Irr}(G)} \dim \Hom_G(\rho^{\otimes k}\otimes \mu, \mu)= tr(A^k_{\Gamma})$$ as required.
\end{proof}

\begin{coro}\label{cor:trace}
 Assume $\Gamma$ is undirected and loop-less. Then $$\dim \Hom_G(\rho^{\otimes 2}, \C[G]^{adj}) =\sum_{g \in G//G} \chi_\rho(g)^2 = 2\abs{\{\text{undirected edges in } \Gamma\}}.$$
\end{coro}




\subsection{Bipartite McKay graphs}
\begin{lemma}\label{lem:bipartite_eigen}
 Let $\Gamma$ be a bipartite graph with adjacency matrix $A_{\Gamma}$. Then for each eigenvalue $\lambda$ of $A_{\Gamma}$, $-\lambda$ is also an eigenvalue of $A_{\Gamma}$.
\end{lemma}
\begin{proof}
 Let $X(\Gamma) = X_0 \sqcup X_1$ be the bipartition of the set of vertices, such that $$\forall i\in \{1,2\}, \; \forall x, y \in X_i, \; x\notin \mathtt{N}(y).$$ Let $f: X(\Gamma) \to \C$ be the eigenvector of $A_{\Gamma}$ with eigenvalue $\lambda$. Then $$\forall x\in X(\Gamma), \; \lambda f(x) = \sum_{y \in \mathtt{N}(x)} f(y).$$ Now define a new function $\tilde{f}: X(\Gamma) \to \C$ by setting $$\forall i\in \{1,2\}, \; \forall x\in X_i, \; \tilde{f}(x) \coloneqq  (-1)^i f(x).$$ Then for any $x\in X_i$, we have: 
 $$-\lambda \tilde{f}(x) = -\lambda (-1)^i f(x) = \sum_{y \in \mathtt{N}(x)}(-1)^{i+1} f(y) =  \sum_{y \in \mathtt{N}(x)} \tilde{f}(y),$$ making $\tilde{f}$ an eigenvector of $A_{\Gamma}$ with eigenvalue $-\lambda$.
\end{proof}

Let $\Gamma \coloneqq  \Gamma(G, \rho)$ be the McKay graph for a finite group $G$ and a representation $\rho$ of~$G$.
\begin{lemma}\label{lem:bipartite_center}

 Assume $\rho$ is irreducible and faithful. Then the graph $\Gamma= \Gamma(G, \rho)$ is bipartite if and only if $C_2 \subset Z(G)$. 
\end{lemma}

\begin{proof}

$[\mathbf{\Leftarrow}]$: Assume $C_2 \subset Z(G)$, and let $z\in C_2\subset Z(G)$, $z\neq 1_G $. Then the vertices in $\Gamma$ can be partitioned into two disjoint subsets, according to the eigenvalue by which $z$ acts on the given irreducible representation of $G$. Since $\rho$ is faithful and irreducible, $\rho(z) =-\id$. So clearly, there are no edges in $\Gamma$ between vertices $\lambda, \mu \in \mathrm{Irr}(G)$ such that $\lambda(z) = \mu(z)$. Hence $\Gamma$ is bipartite.

 $[\mathbf{\Rightarrow}]$: Assume $\Gamma$ is bipartite. By Lemma \ref{lem:bipartite_eigen}, for each eigenvalue $\lambda$ of $A_{\Gamma}$, we also have an eigenvalue~$-\lambda$.  
 
 
 Next, the eigenvalues of $A_{\Gamma}$ are precisely $\chi_\rho(g)$ by Proposition \ref{prop:spectra_adj_matrix}, so there exists $z \in G$ with 
 $$\chi_\rho(z) = -\chi_{\rho}(1_G) = - \dim \rho$$
 Since $z \in G$ has finite order, $\rho(z)$ is a diagonalizable operator whose eigenvalues are all roots of unity. We have $tr(\rho(z)) = - \dim \rho$ and hence $\rho(z) = -\id$. Faithfulness of $\rho$ now implies that $z \in Z(G)$ and $z^2 = 1_G$. Hence $C_2 \subset Z(G)$ as required.
 
\end{proof}

\begin{lemma}\label{lem:faithful_center}
 Assume $\rho$ is irreducible and faithful, and that $\rho^* \cong \rho$. Then $Z(G) \subset C_2 $.
\end{lemma}

\begin{proof}
  By Schur's lemma, $\rho$ defines a character of the center $$\lambda:Z(G) \to \C^{\times}, \;\, z\mapsto \lambda(z),$$ where $\rho(z) = \lambda(z)\id$. Since $\rho^* \cong \rho$, we have: $\triv \subset \rho \otimes \rho^* \cong \rho^{\otimes 2}$ and on this space $z \in Z(G)$ acts by $\lambda^2(z)\id$. So $\lambda(z)^2 = 1$ and thus $\lambda(z)=\pm 1$. Now, since $\rho$ is faithful, the map $\lambda: Z(G) \to \C$ is injective, hence $Z(G) \subset C_2 $.
\end{proof}

\begin{rema}
 {These statements relate to the fact that the group of central characters $U\coloneqq \Hom(Z(G), \C^{\times})$ defines a universal grading on the tensor category $\Rep(G)$, and thus means that $\Gamma(G, \rho)$ is a $\abs{Z(G)}$-partite graph for any $\rho \in \Rep(G)$. Furthermore, if $\rho$ is irreducible and $\zeta_{\rho} \in U$ is the central character corresponding to $\rho$, every closed walk in $\Gamma(G, \rho)$ has length divisible by the order of $\zeta_{\rho}$ in $U$, which implies the statement of Lemma \ref{lem:faithful_center}.}
\end{rema}

Recall that by Lemma \ref{lem:undirected}, $\rho^* \cong \rho$ if and only if $\Gamma$ is undirected, and by Lemma~\ref{lem:connectivity}, $\rho$ is faithful if and only if $\Gamma$ is connected. Using Lemmas \ref{lem:bipartite_center}, \ref{lem:faithful_center}, we conclude:

\begin{coro}\label{cor:bipartite_center_Z_2}
 Assume $\rho$ is irreducible and $\Gamma = \Gamma(G, \rho)$ is an undirected and connected graph. Then $\Gamma$ is bipartite if and only if $Z(G) = C_2 $.
\end{coro}
